%
% lemniskate.tex -- Lemniskate von Bernoulli, r^2 = cos(2 phi)
%
% Normierung a = 1/sqrt(2), Scheitel bei (+-1,0).
% Der Bogen vom Nullpunkt bis zum Scheitel (1,0) ist hervorgehoben,
% er entspricht dem Integral s(1) = varpi/2.
%
\documentclass[tikz]{standalone}
\usepackage{amsmath}
\usepackage{times}
\usepackage{txfonts}
\usetikzlibrary{arrows,math}
%
% farben.tex -- Farben aus buch/common/macros.tex
%
\definecolor{darkgreen}{rgb}{0,0.6,0}
\definecolor{darkred}{rgb}{0.8,0,0}
\definecolor{orange}{rgb}{1,0.6,0}
\definecolor{gelb}{rgb}{1,0.8,0}

\begin{document}
\def\skala{2.6}
\begin{tikzpicture}[>=latex, thick, scale=\skala]

% Achsenlabels verschieben
\def\xdx{-0.05}\def\xdy{-0.05}
\def\ydx{-0.05}\def\ydy{0.0}

% Achsen
\draw[->] (-1.35,0) -- (1.35,0);
\draw[->] (0,-0.65) -- (0,0.65);
\node at ({1.35+\xdx},{\xdy}) [anchor=north east] {$X$};
\node at ({\ydx},{0.5+\ydy}) [anchor=south east] {$Y$};

% rechtes und linkes Blatt
\draw plot[domain=-45:45, samples=100, smooth, variable=\t]
	({sqrt(cos(2*\t))*cos(\t)}, {sqrt(cos(2*\t))*sin(\t)});
\draw plot[domain=135:225, samples=100, smooth, variable=\t]
	({sqrt(cos(2*\t))*cos(\t)}, {sqrt(cos(2*\t))*sin(\t)});

% hervorgehobener Bogen vom Nullpunkt bis zum Scheitel (1,0)
\draw[darkred, line width=1.8pt] plot[domain=0:45, samples=60, smooth, variable=\t]
	({sqrt(cos(2*\t))*cos(\t)}, {sqrt(cos(2*\t))*sin(\t)});

% Radius r bei phi = 25 Grad auf dem hervorgehobenen Bogen
\def\rphi{20}
\draw[thin] (0,0) -- ({sqrt(cos(2*\rphi))*cos(\rphi)}, {sqrt(cos(2*\rphi))*sin(\rphi)})
	node[midway, sloped, above] {$r$};
\fill ({sqrt(cos(2*\rphi))*cos(\rphi)}, {sqrt(cos(2*\rphi))*sin(\rphi)}) circle (0.02);

% Scheitelbeschriftung verschieben
\def\pdx{-0.18}\def\pdy{0.18}
\def\mdx{0.22}\def\mdy{0.18}

% Nullpunkt und Scheitel
\fill (0,0) circle (0.03);
\fill (1,0) circle (0.03);
\fill (-1,0) circle (0.03);
\node[anchor=north] at ({1+\pdx},{\pdy}) {$(1,0)$};
\node[anchor=north] at ({-1+\mdx},{\mdy}) {$(-1,0)$};

\end{tikzpicture}
\end{document}
